\documentclass{article}
\usepackage{graphicx} % Required for inserting images

\usepackage[utf8]{inputenc}
\usepackage[serbian]{babel}
\usepackage[hidelinks]{hyperref}
\usepackage{csquotes}
\usepackage[backend=biber, style=numeric, sorting=none]{biblatex} % sortiranje referenci po redosledu citiranja
\addbibresource{literatura.bib} 

\graphicspath{ {./slike/} }

\usepackage{titlesec}

% Povećavanje sekcija
\titleformat{\section}{\Large\bfseries}{\thesection}{1em}{}
\titleformat{\subsection}{\large\bfseries}{\thesubsection}{1em}{}
\titleformat{\subsubsection}{\normalsize\bfseries}{\thesubsubsection}{1em}{}



\usepackage{geometry}
\geometry{
 a4paper,
 % total={170mm,257mm},
 left=25mm,
 top=25mm,
 right=25mm,
 bottom=25mm
 }
 
\title{Računarstvo i društvo -- Pregled gradiva}
\author{}
\date{}

\begin{document}
\maketitle

\tableofcontents
\newpage


\section{Istorijski pogled razvoja računara}

\subsection{Pregled gradiva}
Razvoj računara otpočinje pojavom mehaničkih uređaja za obradu podataka. Poseban značaj imaju bušene kartice Hermana Holerita, koje su znatno ubrzale obradu popisa stanovništva u SAD krajem XIX veka. Iz ovog tehnološkog rešenja nastaju velike kompanije poput IBM-a koje su oblikovale korene industrije računara.

Sredinom XX veka počinje razvoj elektronskih računara, od Atanasof-Berijevog ABC računara, preko ENIAC računara i fon Nojmanove arhitekture, pa sve do prvih komercijalnih sistema kao što su UNIVAC i IBM PC. Razvojem tranzistora i integrisanih kola računari postaju manji, jeftiniji i dostupniji. Ova promena dovodi do razvoja kompanija Microsoft i Apple i  pojave personalnih računara.

Paralelno se radi i na komunikaciji, od telegrafa i telefona, preko radija i televizije, do interneta. Pojavljuje se i ARPANET tehnologija koja omogućava povezivanje računara radi razmene podataka u paketima. ARPANET se smatra pretečom interneta. Kasnije se razvijaju i TCP/IP protokol, mejlovi, World Wide Web koji kao rezultat daju globalnu digitalnu povezanost. Grafički korisnički interfejs i veb pretraživači podižu tehnologiju na nivo koji je dostupan svima.


\subsection{Predlozi tema}
\begin{enumerate}
    \item Doprinos žena razvoju računarsva i postavljanju temelja programiranja u XX veku
        \begin{description}
            \item Ada Lovelace kao prva programerka.  \href{https://www.youtube.com/watch?v=IZptxisyVqQ}{link} 
            \item Žene programeri kompanije ENIAC. \href{https://www.youtube.com/watch?v=Zevt2blQyVs}{link}
        \end{description}
    \item Računari u Jugoslaviji: zaboravljena istorija CER i TARA računara
        \begin{description}
            \item \href{https://www.youtube.com/watch?v=T8iebSL49NM}{link}
            \item \href{https://www.rts.rs/lat/magazin/tehnologija/5583359/dan-informaticara-srbije-prvi-racunar-cer10-18-novembar-1960.html}{link}
        \end{description}
    \item Društveni uticaj računara: od prvog popisa do digitalnog doba
        \begin{description}
            \item \href{https://www.youtube.com/watch?v=WWzaZWN1IOU}{link}
        \end{description}
\end{enumerate}


\section{Etika}

\subsection{Pregled gradiva}
Etika je grana filozofije koja proučava moral, sistem pravila koje društvo propisuje kako bi štitilo i unapređivalo dobrobit svojih članova. Moral čine pravila i norme, dok etika nastoji da ih razume.

Etika se sastoji od nekoliko podoblasti:
\begin{enumerate}
    \item Meta-etika - bavi se pitanjima o značenju moralnih pojmova i prirodi moralne obaveze
    \item Normativna etika - postavlja opšta načela i principe
    \item Primenjena etika - razmatra konkretne situacije i dileme
\end{enumerate}

Razmatramo devet etičkih teorija:
\begin{enumerate}
    \item Subjektivni realizam - svaki pojedinac sam za sebe odlučuje šta je dobro, a šta loše
    \item Kulturološki relativizam - značenje dobrog i lošeg zavisi od principa koje definiše konkretno društvo
    \item Teorija duhovnog komandovanja - nešto je dobro ako je u skladu sa Božijom voljom, a loše ako nije
    \item Etički egoizam - moralno dobra akcija je ona koja će toj osobi dugoročno gledano doneti maksimalnu dobit
    \item Kantijanizam - ljudsko delovanje je vođeno moralnim pravilima koja su univerzalna za sve
    \item Utilitarizam postupka - akcija je dobra do mere u kojoj ona koristi svim uključenim stranama
    \item Utilitarizam pravila - prihvata ona moralna pravila koja kada bi ih svi pratili povećavaju ukupnu sreću svih uključenih strana
    \item Teorija socijalnog ugovora - moralnost je skup pravila koja određuju kako se ljudi ponašaju jedni prema drugima koje racionalni ljudi prihvataju zarad opšteg dobrobita pod uslovima da i drugi ta pravila poštuju
    \item Etika vrline - moralno ispravna akcija je ona koju bi u istim okolnostima izabrala i uradila moralna osoba
\end{enumerate}

Razvoj tehnologije stvara nove etičke probleme i kreira stavove koje možemo podeliti na:
\begin{enumerate}
    \item Tehnicizam - tehnologija je korisan alat za rešavanje problema velikog potencijala
    \item Optimizam - tehnologija ima pozitivan uticaj na drušvo
    \item Skepticizam - tehnologija ima negativan uticaj na društvo
\end{enumerate}

Postoje i različiti pristupi računarkoj etici koji se mogu podeliti na:
\begin{enumerate}
    \item Pristup nepostojanja - računarska etika ne postoji kao posebna oblast
    \item Profesionalni pristup - IT stručnjaci treba da imaju poseban profesionalni kodeks ponašanja
    \item Radikalni pristup - razvoj tehnologije je stvorio potpuno nove moralne situacije koje klasična etika ne može da objasni
    \item Konzervativni pristup - razvoj tehnologije nije promenio samu prirodu etičkih problema, već samo kontekst u kojem se oni javljaju
    \item Inovativni pristup - saglasan je da razvoj tehnologije donosi nove izazove i moralne dileme, ali ne odbacuje potpuno postojeću etiku
\end{enumerate}

\subsection{Predložene teme}
\begin{enumerate}
    \item Nadzor i prepoznavanje lica: privatnost, ljudska prava i društvene posledice
        \begin{description}
            \item \href{https://www.youtube.com/watch?v=2SdpzTZTznw}{link}
            \item \href{https://www.youtube.com/watch?v=CK-mkq0HQ2w}{link} 
        \end{description}
    \item Etika lažnih vesti i dezinformacija na internetu
        \begin{description}
            \item \href{https://journals.publishing.umich.edu/jpe/article/id/6195/}{link}
            \item \href{https://www.youtube.com/watch?v=uCESQTkxuN0}{link}
        \end{description}
    \item AI-generisana umetnost i kreativnost
        \begin{description}
            \item \href{https://link.springer.com/article/10.1007/s43681-025-00735-3}{link}
            \item \href{https://ai-arts.org/}{link}
            \item \href{https://aifcf.org/}{link}
        \end{description}
\end{enumerate}

\section{Upotreba interneta, spam poruke}
\subsection{Pregled gradiva}
U današnjem svetu sve više ljudi ima mobilni telefon i pristup internetu. Internet mreže omogućavaju bržu komunikaciju i brže obacljanje svakodnevnih poslova. Ipak, uz sve prednosti, moderna komunikacija donosi i niz etičkih izazova i zloupotreba.

Spam poruke su neželjene poruke, koje stvaraju ogromne finansijske i resursne posledice.
Spam je izuzetno jeftin.  Iako odziv iznosi samo jednu poruku na 12,5 miliona, globalno se godišnje zaradi oko 3,5 miliona dolara, što je dovoljno da spam i dalje opstaje. Najčešće se radi o reklamama, eksplicitnom sadržaju, finansijskim ponudama i prevarama. Spam utiče na ekonomiju firmi i funkcionisanje infrastrukture: zaposleni gube vreme na brisanje poruka, što otežava rad, a Internet i serverski resursi bivaju oštećeni.

Tehnologija se razvija kako bi suzbila spam. Koriste se spam-filteri na nivou provajdera i korisnika, Bayesova filtracija, crne liste servera, CAPTCHA-izazovi, challenge-response sistemi i verifikacija adrese pošiljaoca.

Internet i World Wide Web (WWW) nisu isto. Internet je globalna mreža računara, a WWW je skup međusobno povezanih dokumenata na toj mreži. Web pregledači i mobilne aplikacije omogućavaju pristup toj mreži.
Internet ne služi samo zabavi, već i za online kupovinu, društvene mreže, deljenje sadržaja, obrazovanje, ali i internet klađenju i kockanju i drugom sadržaju koji otvara etička pitanja o dostupnosti, zaštiti dece i regulaciji.

\subsection{Predložene teme}
\begin{enumerate}
    \item Spam i ekološki uticaj digitalnog otpada
        \begin{description}
            \item \href{https://www.theguardian.com/environment/green-living-blog/2010/oct/21/carbon-footprint-email}{link} 
            \item \href{https://blog.cabi.org/2011/08/05/ever-wondered-what-the-carbon-footprint-of-your-spam-emails-is/}{link}
        \end{description}
    \item Etika automatizacije i botova u slanju poruka
        \begin{description}
            \item \href{https://www.forbes.com/councils/forbescoachescouncil/2024/11/06/ethical-ai-in-marketing-balancing-automation-with-human-values/}{link}
            \item \href{https://mint.rs/blog/bezbednost/zasto-e-mail-poruke-zavrsavaju-u-spam-folderu/}{link}
        \end{description}
\end{enumerate}

\section{Cenzura i sloboda govora, deca i neprikladan sadržaj}
\subsection{Pregled gradiva}
Cenzura je pokušaj regulacije ili zabrane javnog pristupa informacijama i materijalima koji se smatraju štetnim, opasnim ili uvredljivim.
Postoje dve osnovne forme:
\begin{enumerate}
  \item Direktna cenzura - nametnuta spolja, od strane države ili društva — poput državnog monopola nad medijima, pregleda publikacija ili sistema licenciranja.
  \item Autocenzura - samonametnuto odbacivanje objavljivanja određenih sadržaja često radi sopstvene zaštite ili iz straha da ne povredimo druge.
\end{enumerate}

U današnjem svetu direktna cenzura je i dalje aktivna i ima tri forme:
\begin{enumerate}
  \item Monopol nad informacijama od strane slade ili države
  \item Pregled publikacije
  \item Licenciranje i registracija
\end{enumerate}


Internet se često koristi za širenje dezinformacija i manipulacija. Kao odgovor, vlasti i platforme mogu primeniti vidove cenzure, ne uvek u potpunoj zabrani, već u smanjenom vidljivosti određenog sadržaja. Postoje situacije kada je regulacija opravdana, kao u slučaju zaštite dece od neprikladnog sadržaja. Međutim, granica između zaštite i preterane kontrole je često tanka i zahteva pažljivu etičku procenu.


Etika i zakon nisu isto. Zakon sprovode države, etička pravila proizilaze iz društvenih normi. Sa razvojem digitalnih tehnologija, etička pitanja postaju sve kompleksnija. Treba jasno postaviti granice između zaštite i cenzure, uzimajući u obzir ko ima moć, koji je cilj regulacije i koje su posledice po slobodu izražavanja i društvo u celini.

Na internetu je lako dostupan pornografski i nasilan sadržaj. Kako bi se deca zastitila od neprikladnog sadržaja dizajnirani su veb filteri. Veb filteri su programi koji sprečavaju određenim Veb stranicama da budu prikazani u Veb pregledaču. Često su korišćeni od strane roditelja, kompanija ili obrazovnih institucija.
Još jedan problem dece ili adolescenata na internetu jeste seksting. Seksting predstavlja slanje poruka ili imejlova koje imaju seksualno sugestivni karakter ili slanje slika na kojima su osobe sa vrlo malo odeće ili bez odeće.

\subsection{Predložene teme}
\begin{enumerate}
    \item Digitalni trag i privatnost korisnika
        \begin{description}
            \item \href{https://www.youtube.com/watch?v=cprsK9uoYjk&t=2s}{link}
            \item \href{https://www.internetsociety.org/blog/2024/10/understanding-digital-footprints/}{link}
        \end{description}
    \item Cyberbullying i digitalna odgovornost
        \begin{description}
            \item \href{https://cyberpsychology.eu/article/view/35784}{link}
            \item \href{https://www.frontiersin.org/journals/psychology/articles/10.3389/fpsyg.2022.861823/full}{link}
        \end{description}
\end{enumerate}

\section{Internet prevare i zavisnost od interneta}
\subsection{Pregled gradiva}
Razvoj računara i interneta doneo je ogromne mogućnosti, ali i ozbiljne rizike. Jedan od značajnijih problema su prevare i širenje netačnih ili fabrikovanih informacija koje se šire kroz digitalne kanale.

Krađa identiteta predstavlja ozbiljnu pretnju. Počinilac koristi tuđe ime, dokumente, matični broj ili podatke sa bankovnih kartica kako bi ostvario prava i transakcije koje pripadaju žrtvi.
Phishing, pecanje, je najčešći oblik prevare na internetu. Napadači šalju poruke koje izgledaju kao da potiču od poznatih servisa poput Amazon-a, PayPal-a, eBay-a i traže od žrtve da unese lične podatke kao korisničko ime, lozinku, podatke o bankovnim računima i slično.
Prevaranti često koriste lažne Veb stranice koje oponašaju legitimne, ali nemaju bezbednosne oznake poput HTTPS sertifikata. Povećana upotreba mobilnih telefona dodatno otežava prepoznavanje prevara.
Nigerijska prevara, poznata još iz 1980-ih, bazirana je na lažnim pričama gde prevarant tvrdi da ne može prebaciti novac zato što nailazi na administrativne probleme, moli žrtvu za pomoć uz obećanje o velikoj zaradi.

Još jedan vid internet prevara su i netačne informacije. Internet je preplavljen informacijama, kako korisnim i edukativnim tako i onim netačnim. Postoje različiti mehanizmi sortiranja stranica na internetu među kojima je Google ocenjivanje kvaliteta stranica sistemom glasanja.

Među najvećim problemima interneta ubraja se i zavisnost od interneta. U nekim zemljama se ova zavisnost posmatra kao psihička bolest. Ova zavisnost izaziva probleme poput depresije, anksioznosti, promene raspoloženja, usameljenosti, glavobolje, insomnije i slično.


\subsection{Predložene teme}
\begin{enumerate}
    \item Psihologija phishing-a: zašto ljudi klikću i kad znaju da ne bi smeli
        \begin{description}
            \item \href{https://www.sciencedirect.com/science/article/pii/S0167404823005813}{link} 
            \item \href{https://www.youtube.com/watch?v=AQLxyjuDTwQ}{link} 
        \end{description}
    \item Digitalna detoksikacija: lečenje zavisnosti od interneta
        \begin{description}
            \item \href{https://www.youtube.com/watch?v=SKNjQtkeJWg}{link}
            \item \href{https://www.verywellmind.com/internet-addiction-4157289}{link}
        \end{description}
    \item Digitalna hiperpažnja: deca i vizuelni stimulansi 
        \begin{description}
            \item \href{https://www.frontiersin.org/journals/psychology/articles/10.3389/fpsyg.2020.00961/full}{link}
            \item \href{https://www.sciencedirect.com/science/article/pii/S0360131521001332}{link}
            \item \href{https://www.youtube.com/watch?v=3h1T493piLE}{link} 
        \end{description}
\end{enumerate}

\section{Bezbednost interneta}

\subsection{Pregled gradiva}
Značenje reči haker se menjalo vremenom. U prošlom veku haker je označavao istraživača, osobu koja preuzima rizik, nekoga ko čini da sistem radi nešto novo. Danas se termin haker se vezuje za osobu koja neovlašćeno pristupa računarima i računarskim mrežama.
Među korišćenim tehnikama hakovanja možemo izdvojiti:
\begin{itemize}
    \item Metod grube sile - za kratke šifre
    \item Metod rečnika - za duge šifre
\end{itemize}

Hakovanje u okviru besplatne WiFi mreže je vrsta napada u kome haker pristupi kolačiću koji stranica šalje korisniku. Na osnovu kolačića, haker dobija pristup sajtu i ista prava i mogućnosti kao i registrovani korisnik koji je napadnut. Prema Krivičnom zakoniku Republike Srbije kazne za hakere su novčane ili zatvorske, u zavisnosti od težine krivičnog dela.
Zlonamerni softver je softver dizajniran da nanese štetu ciljnom korisniku. Zlonamerne softvere delimo u više grupa:
\begin{itemize}
    \item Virusi - programski kod koji se samoreplicira i deo je drugog programa
    \item Računarski crvi
        \begin{itemize}
            \item Internet crvi - samostalni program koji se širi računarima u mreži rupa u bezbednosti mreže
            \item Saser - program koji je isključivao računare neko vreme nakon podizanja sistema Windows XP ili Windows 2000
        \end{itemize}
    \item Ukrštanje veb lokacija - na stranicama gde korisnici mogu ostaviti komentare, zlonamerni korisnici prave skriptu koja drugim korisnicima prati kolačiće
    \item Usputno preuzimanje podataka - napad na računar posetioca zaražene veb stranice, tako što se na njegov računar kopiraju zlonamerne datoteke
    \item Trojanski konj - štetan program, koji je maskiran u program koji je korisan, dok u pozadini izvršava zlonamerne akcije bez znanja korisnika
    \item Softver za špijuniranje - neželjeni softver koji upada u sistem i krade poverljive informacije
    \item Mreža botova - kolekcija međusobno povezanih zaraženih uređaja, koji su pod kontrolom istog tipa zlonamernog softvera
\end{itemize}

Postoji i lista mera koje je poželjno primenjivati kako bi se izbegle gore navedene neprijatnosti. 
\begin{itemize}
    \item Bezbednosne zakrpe softvera - ažurirane delove softverskog rešenja koje je u svim prethodnim verzijama imalo bezbedonosne propuste
    \item Alati za odbranu od zlonamernog softvera - koriste se u preventivne svrhe ili za analizu i otklanjanje zaraženih datoteka
    \item Zaštitni zidovi - softver za regulisanje mrežnog saobraćaja
\end{itemize}
\subsection{Predložene teme}
\begin{enumerate}
    \item Evolucija zlonamerkog softvera 
        \begin{description}
            \item od prvih virusa do AI napada
            \item \href{https://www.kaspersky.com/resource-center/threats/a-brief-history-of-computer-viruses-and-what-the-future-holds}{link} 
            \item \href{https://www.forbes.com/councils/forbestechcouncil/2024/03/19/how-ai-driven-cyber-attacks-will-reshape-cyber-protection/}{link}
        \end{description}
    \item „Pametni dom“ kao meta hakera
        \begin{description}
            \item \href{https://www.quectel.com/}{link} 
            \item \href{https://www.cisco.com/site/us/en/products/security/industrial-security/index.html}{link}
        \end{description}
\end{enumerate}

\section{Privatnost na internetu}
\subsection{Pregled gradiva}

Privatnost je često teško definisati jer obuhvata zaštitu od neovlašćenog nadzora, kontrolu nad sopstvenim telom i informacijama, kao i pravo na slobodu mišljenja i govora. Ključan pojam je „pristup“, odnosno pravo pojedinca da ograniči ko i koliko vidi njegove podatke, ali i balans s pravom drugih da pristupe određenim informacijama.
Kada govorimo o privatnosti na društvenim mrežama bitno je naglasiti da bez obzira na to koliko neko odluči da deli informacije, to što njegovi prijatelji odluče da dele takođe igra značajnu ulogu, čak i ako pojedinac nema nalog na društvenim mrežama.

Kolčići su male tekstualne datoteke koje sajtovi koriste za pamćenje informacija poput prijava, sadržaja korpe i slično. Postoje sesijski (kratkotrajni) i trajni kolčići, koji imaju rok trajanja i čuvaju se na računaru. Trajni kolačići se dele na kolačiće prve strane, koji se priste samo od strane veb sajtova koji su ih napravili, i na kolačiće treće strane koji prate korisnika na drugim domenima
Zahvaljujući zakonima, sajtovi moraju da traže direktnu saglasnost za upotrebu kolčića, što korisnici često ne razumeju jer ne čitaju pravila privatnosti.
Države koriste nadzor kako bi očuvale sigurnost, ali često prelaze etičke i zakonske granice.

Neki od glavnih tehnoloških mehanizama zaštite su:
\begin{itemize}
  \item Enkripcija -  sprečava da podaci budu čitljivi u tranzitu
  \item Virtualne privatne mreže -  sakrivaju adrese i štite komunikaciju
\end{itemize}

\subsection{Predložene teme}
\begin{itemize}
  \item Big Data i etika privatnosti
    \begin{description}
        \item \href{https://hbr.org/2023/07/the-ethics-of-managing-peoples-data}{link} 
        \item \href{https://www.youtube.com/watch?v=bAyrObl7TYE}{link}
    \end{description}
\end{itemize}


\end{document}
